%% ============================================================
\section{Related Work}\label{sec:related}
%% ============================================================

\paragraph{Classical invariant generation.}
Static approaches compute invariants as fixpoints in abstract domains: Karr's affine-relation analysis~\citep{karr1976affine}, abstract interpretation~\citep{cousot1977abstract}, and the polyhedral domain~\citep{cousot1978automatic} discover invariants of a fixed shape without any target, but are limited to the chosen domain's expressiveness and lose precision under widening.
Constraint- and template-based methods fix a syntactic template and solve for coefficients---linear invariants via Farkas-style non-linear constraint solving~\citep{colon2003linear}, non-linear invariants via Gr\"obner bases~\citep{sankaranarayanan2004non}, and program analysis recast as constraint solving~\citep{gulwani2008program,gupta2009invgen}.
In verification-condition frameworks, invariants are by-products of proving a property: predicate abstraction~\citep{graf1997predicate,ball2002slam}, Craig interpolation~\citep{mcmillan2003interpolation}, IC3/PDR~\citep{bradley2011sat,komuravelli2014smt,gurfinkel2015seahorn}, and abduction~\citep{dillig2013inductive} all generalize from the \emph{negation of the target}; without a property to refute, these engines have nothing to generalize from.
Houdini~\citep{flanagan2001houdini} computes the greatest inductive subset of a candidate pool; \sam{} reuses exactly this algorithm as its soundness backstop, but delegates \emph{candidate generation}---the part Houdini leaves to heuristic annotation guessing---to a learned policy.
Furia and Meyer generate invariant candidates by syntactically mutating the postcondition~\citep{furia2010inferring}, the clearest evidence for our motivating observation that in practice the target and the invariant are two projections of the same object.

\paragraph{Data-driven and learning-based invariant inference.}
Daikon pioneered inferring \emph{likely} invariants from execution traces without any specification~\citep{daikon}; DIG extended this to polynomial and array invariants~\citep{nguyen2012using}, later with symbolic-state checking~\citep{nla2019,nguyen2017syminfer}.
These works share our target-free spirit but stop at ``likely'': candidates are unsound until confirmed, and confirmation is left to a downstream prover.
Guess-and-check frameworks close that loop: invariant inference as classification over reachable/unreachable states~\citep{sharma2012interpolants,sharma2013data}, ICE learning with implication counterexamples~\citep{garg2014ice}, its decision-tree instantiation~\citep{garg2016learning}, Horn-ICE for contracts~\citep{ezudheen2018horn}, data-driven CHC solving~\citep{zhu2018data}, and PIE's on-demand feature learning~\citep{padhi2016pie}.
Crucially, in all of these the \emph{negative} and \emph{implication} examples are produced by a verifier attempting a given property: the teacher is target-directed even when the learner is agnostic.
\sam{} keeps the classification view---an invariant is scored by how it separates reachable from unreachable states---but manufactures negatives \emph{from the loop's own dynamics} (impossible continuations and off-manifold perturbations of real traces), so the signal exists before any property is stated.

\paragraph{Neural and LLM-based invariant generation.}
Code2Inv learns a policy over invariant syntax trees with reinforcement from solver feedback~\citep{si2018learning,si2020code2inv}; CLN2INV and G-CLN learn invariant coefficients from traces with differentiable logic~\citep{ryan2020cln2inv,yao2020gcln}.
With LLMs, \citet{pei2023can} fine-tune to predict Daikon-style invariants at scale; Loopy prompts GPT-4 for ACSL invariants repaired by Houdini~\citep{kamath2023finding}; iRank re-ranks LLM candidates~\citep{chakraborty2023ranking}; Lemur interleaves LLM proposals with automated reasoners~\citep{wu2024lemur}; AutoSpec and SpecGen synthesize full specifications with static-analysis guidance~\citep{wen2024enchanting,ma2024specgen}; benchmarks now probe memory-manipulating programs~\citep{liu2024towards} and proof-oriented languages~\citep{chakraborty2024towards,first2023baldur,misu2024towards}.
Almost without exception these systems receive the assertion in the prompt and are judged---and, where trained, rewarded---by whether the target verifies.
This yields exactly the failure mode we target: sparse, binary, hackable feedback~\citep{skalse2022defining,gao2023scaling} that favors restating $Q$ over understanding the loop.
\sam{} differs on both axes: the generator is \emph{closed-book} (the postcondition is stripped from its input), and the reward is a dense, execution-grounded measure of summary tightness rather than a verifier bit.

\paragraph{Loop summarization.}
Summarizing loop effects independently of any property has a long history in static analysis: abstract-transformer summaries~\citep{kroening2008loop}, partial summaries in dynamic test generation~\citep{godefroid2011automatic}, disjunctive summaries via path dependency~\citep{xie2016proteus}, compositional recurrence analysis~\citep{farzan2015compositional,kincaid2017compositional}, and acceleration~\citep{bardin2008fast}.
These techniques are precise on the loop classes their theories cover but brittle outside them.
We adopt the task framing---characterize the loop, not the goal---but express summaries as ACSL inductive invariants checked by Frama-C/WP~\citep{framac,acsl}, and use learning to cover loop shapes (nondeterministic guards, branch-dependent updates, non-linear laws) that closed-form summarization does not.

\paragraph{RL from verifiable rewards and reward hacking.}
Post-training LLMs against programmatically checkable signals---unit tests for code~\citep{le2022coderl,gehring2024rlef,chen2021evaluating}, final answers for mathematics~\citep{shao2024deepseekmath,lambert2024tulu,deepseekr1}, or process supervision~\citep{lightman2023lets}---has proven far more robust than learned preference models~\citep{ouyang2022training,rafailov2023direct}, yet even verifiable rewards are gamed when the checker under-specifies the intent~\citep{amodei2016concrete,skalse2022defining,gao2023scaling}.
Our setting sharpens this problem: the reward is computed from a \emph{finite sample} of program behavior, so memorizing the sample is always a candidate policy.
The layered defenses of $\S$\ref{sec:method:antihack}---capacity limits on the policy's output, uniqueness-based credit, and minimum-scoring across decorrelated held-out samples---are, to our knowledge, the first systematic anti-hacking treatment for an execution-based formal-methods reward, and we validate them with adversarial gates rather than by assumption.
